\section{Confiabilidad, Plan Divergence y LLM OS}

Al final, el curso concentra los problemas pendientes del agent en la reliability. El software tradicional puede ejecutarse repetidamente con la misma entrada; la agent policy, en cambio, interactúa con un entorno dinámico a lo largo de muchos stochastic step. Aunque la tasa de éxito local parezca alta, en las trayectorias largas los fallos se acumulan con rapidez.

\subsection{Problemas centrales de los agents autónomos}

La diapositiva Key Issues enumera reliability, looping, plan divergence, benchmark, observabilidad, security y human fallback. No son funciones complementarias que se añaden después del despliegue, sino parte de la agent architecture.

\lecturefigure{slide-109.jpg}{Problemas de confiabilidad, ciclos, pruebas y seguridad de un autonomous agent}{Slide deck oficial de CS25 V4, página 109}

\begin{importantbox}{Efecto multiplicativo en la tasa de éxito de trayectorias largas}
Si cada paso tiene una probabilidad de éxito independiente $p$, la probabilidad de que una trajectory de longitud $T$ tenga éxito completo es aproximadamente
\[
P(\text{all success})=p^T.
\]
Incluso con $p=0.99$, al ejecutar 100 pasos la probabilidad de terminar sin ningún error es de apenas $0.99^{100}\approx 36.6\%$. Los errores reales no son independientes, y una desviación temprana además modifica las observation posteriores, lo que agrava el fallo.
\end{importantbox}

\teachervoice{El ponente usa la frase «aun con un agent al 95\%, el 5\% restante va a fallar» para explicar la diferencia en el despliegue de software: en un programa tradicional se puede corregir un deterministic bug, mientras que los pocos fallos de un stochastic agent reaparecen una y otra vez a escala de muchos usuarios y en trayectorias largas.}

\begin{knowledgebox}{Qué debe registrar la observabilidad}
\begin{enumerate}
  \item La observation, la selected action y la tool response de cada paso;
  \item Las versiones de model/prompt/tool y el memory snapshot;
  \item La permission decision, la user confirmation y el side effect;
  \item El verifier verdict, los retry, el fallback y el final outcome.
\end{enumerate}
El propósito de los registros es permitir la reproducción y la localización de fallos, no conservar contenido sensible de forma ilimitada; deben acompañarse de minimización, cifrado y una retention policy.
\end{knowledgebox}

\subsection{Plan divergence: cómo volver al rumbo después de desviarse}

La diapositiva de plan divergence contrasta el ideal path con el actual path. Si en algún paso el agent interpreta mal una página, una herramienta devuelve un error o la descomposición del objetivo es incorrecta, las action posteriores pueden apoyarse en un state equivocado y terminar en un loop. Esta sección retoma la reliability de la diapositiva anterior y explica con más detalle por qué el sistema debe comparar periódicamente el state previsto en el plan con la observation real.

\lecturefigure{slide-110.jpg}{Plan divergence: la trayectoria real se aleja paso a paso de la ruta ideal}{Slide deck oficial de CS25 V4, página 110}

\begin{importantbox}{Recovery policy}
El agent no debe preguntarse solo «cuál es el siguiente paso», sino también verificar periódicamente
\[
d_t=D(s_t,\hat{s}_t,g),
\]
donde $s_t$ es la observation real, $\hat{s}_t$ es el state esperado según el plan y $g$ es el objetivo. Cuando el divergence score $d_t$ supera un umbral, el sistema debe pausar, volver a observar, revertir a un checkpoint, replanificar o solicitar intervención humana.
\end{importantbox}

\begin{warningbox}{Un loop detector no puede limitarse a contar texto repetido}
El agent puede repetir la misma action fallida con otras palabras, o bien quedar en un ciclo entre páginas. La detección debe combinar action signature, state hash, error pattern, progress metric y consumo de presupuesto. Al alcanzar el límite de pasos, tiempo o costo, es obligatorio detenerlo de forma definitiva.
\end{warningbox}

\teachervoice{El ponente califica a AutoGPT como un buen prototype que, sin embargo, «en la práctica no puede hacer gran cosa», porque tiende a entrar en ciclos, desviarse y cometer errores de manera continua. Lo importante aquí no es descalificar el prototipo, sino señalar que un sistema de trayectorias largas sin error correction no puede recuperarse por sí solo con solo seguir generando.}

\subsection{LLM OS: una analogía de sistema útil pero incompleta}

El diagrama LLM OS de Karpathy trata al LLM como la CPU, a la context window como la RAM, al embedding store como el file system, a calculator, Python y terminal como software tradicional o ALU, y a las imágenes, el audio, el navegador y otros modelos como peripheral.

\lecturefigure{slide-111.jpg}{Correspondencia de componentes del LLM OS de Karpathy}{Slide deck oficial de CS25 V4, página 111}

\begin{knowledgebox}{Lectura de la figura: tabla de correspondencias del LLM OS}
\begin{tabularx}{\textwidth}{L{0.22\textwidth} L{0.24\textwidth} X}
\toprule
Sistema tradicional & Analogía en el LLM system & Garantías que aún faltan \\
\midrule
CPU & Foundation model & Ejecución determinista, instrucciones tipadas y corrección. \\
RAM & Context window & Aislamiento, consistencia de paginación y escritura controlada. \\
File system & Retrieval/memory store & Transacciones, permisos, versiones y semántica de borrado. \\
ALU/software & Calculator, Python, tools & Sandbox, límites de recursos y semántica de errores. \\
Peripheral/network & Vision, audio, browser, other agents & Autenticación, fronteras de confianza y validación de entradas. \\
\bottomrule
\end{tabularx}
\end{knowledgebox}

\begin{warningbox}{La palabra «OS» no aporta por sí sola las capacidades de un sistema operativo}
Un OS real se encarga de scheduling, memory protection, process isolation, permission, interrupt y recovery. Si un LLM orchestration framework carece de estos mecanismos, solo tiene la forma de un OS, sin sus garantías de seguridad y confiabilidad.
\end{warningbox}

\subsection{Generalized AI system y action engine}

La siguiente diapositiva conecta la chat interface del usuario, el action engine, los agents, la memory, las tools y el mundo exterior en un generalized system. El action engine se encarga de route, plan, permission y aggregation, y es el control plane más importante fuera del modelo.

\lecturefigure{slide-112.jpg}{Diagrama del sistema con chat interface, action engine, agents, tools y memory}{Slide deck oficial de CS25 V4, página 112}

\teachervoice{La visión final sitúa al action engine, detrás de la chat interface, como el centro de enrutamiento: reparte las tareas entre distintos agents y luego integra los resultados. Esa posición también debe asumir permisos, presupuesto, verificación y cancelación, y no puede ser solo un prompt router.}

\begin{importantbox}{Control plane y data plane}
\begin{itemize}
  \item \textbf{Control plane}: interpreta el objetivo, selecciona agent/tool, asigna permisos, mantiene el state y decide retry/fallback;
  \item \textbf{Data plane}: ejecuta efectivamente las action de search, code, browser, database, message o device.
\end{itemize}
Al separar ambos, se puede impedir que una model proposal se convierta directamente en un side effect real.
\end{importantbox}

\subsection{Future needs: error correction, permission y sandbox}

La última diapositiva no concluye que «los agents ya están resueltos», sino que enumera la infraestructura pendiente: error correction, better framework, security, user permission, risky-domain stability y sandboxing. Al leer la figura, conviene tratar estos puntos como deployment precondition y relacionarlos uno por uno con la long-horizon failure, los permisos y los side effect vistos antes.

\lecturefigure{slide-113.jpg}{Los AI agents aún requieren error correction, security, permission y sandbox}{Slide deck oficial de CS25 V4, página 113}

\begin{importantbox}{Pila mínima de seguridad para action de alto riesgo}
\begin{enumerate}
  \item \textbf{Least privilege}: otorgar solo la credential y el action scope mínimos necesarios para la tarea actual;
  \item \textbf{Confirmation gate}: los pagos, envíos, eliminaciones, publicaciones y cambios de permisos requieren confirmación explícita;
  \item \textbf{Sandbox}: el código, los archivos y el acceso a la red se ejecutan en un entorno aislado;
  \item \textbf{Budget}: limitar tokens, pasos, tiempo, costo y llamadas externas;
  \item \textbf{Audit and recovery}: registrar los efectos secundarios y permitir la cancelación, la compensación o la toma de control humana.
\end{enumerate}
\end{importantbox}

\teachervoice{La conclusión final del curso no es que «los agents pronto serán completamente automáticos», sino que la error correction, la security, la user permission y el sandbox siguen sin resolverse. Sobre todo en escenarios de alto riesgo como finance o legal, la estabilidad y la seguridad deben anteponerse a una autonomy más amplia.}

\subsection{Resumen de la sección}

La agent reliability está dominada por el efecto multiplicativo de las trayectorias largas y por la plan divergence. El LLM OS ofrece un lenguaje para organizar los componentes, pero no produce por sí solo isolation, transaction ni recovery. Un sistema desplegable necesita que control plane, least privilege, verification, budget, fallback y sandbox restrinjan en conjunto la model proposal.
